% talk-common.tex —— 各场分享共用的导言区配置
% 放在仓库根目录，由各个入口 .tex 通过 % talk-common.tex —— 各场分享共用的导言区配置
% 放在仓库根目录，由各个入口 .tex 通过 % talk-common.tex —— 各场分享共用的导言区配置
% 放在仓库根目录，由各个入口 .tex 通过 % talk-common.tex —— 各场分享共用的导言区配置
% 放在仓库根目录，由各个入口 .tex 通过 \input{talk-common.tex} 引用。
% 编译时需要把仓库根目录加入 TEXINPUTS，详见根目录 AGENTS.md 第 2 节。

\usetheme[customfont,showmaxslides,nofooterlogo]{pureminimalistic}
\usefonttheme{professionalfonts}
\usepackage[UTF8,fontset=none]{ctex}
\usepackage{graphicx}
\usepackage{listings}
\usepackage{appendixnumberbeamer}

\setsansfont{Fira Sans}
\setmonofont{Fira Mono}
\setCJKmainfont{Noto Sans CJK SC}
\setCJKsansfont{Noto Sans CJK SC}
\setCJKmonofont{Noto Sans CJK SC}

% 默认不引用上游示例 logo，有真实素材时再统一配置
\renewcommand{\logotitle}{}
\renewcommand{\logoheader}{}
\renewcommand{\logofooter}{}
\renewcommand{\pageword}{}

% 字号设置必须放在加载主题之后，避免被主题覆盖
\setbeamerfont{title}{size=\LARGE,shape=\upshape,series=\bfseries}
\setbeamerfont{presentation title}{shape=\upshape}
\setbeamerfont{frametitle}{size=\Large,shape=\upshape,series=\bfseries}

% 对浅色/深色都保持明确的代码对比度
\definecolor{TalkCodeBg}{HTML}{F3F4F6}
\definecolor{TalkCodeFg}{HTML}{17212B}
\definecolor{TalkCodeKeyword}{HTML}{174EA6}
\definecolor{TalkCodeComment}{HTML}{47634D}
\definecolor{TalkCodeString}{HTML}{8A3014}
\lstdefinestyle{talkcode}{
  basicstyle=\ttfamily\small\color{TalkCodeFg},
  backgroundcolor=\color{TalkCodeBg},
  keywordstyle=\bfseries\color{TalkCodeKeyword},
  commentstyle=\color{TalkCodeComment},
  stringstyle=\color{TalkCodeString},
  columns=fullflexible,
  keepspaces=true,
  showstringspaces=false,
  breaklines=true,
  tabsize=4,
  numbers=none,
  frame=none
}
\lstset{style=talkcode}
 引用。
% 编译时需要把仓库根目录加入 TEXINPUTS，详见根目录 AGENTS.md 第 2 节。

\usetheme[customfont,showmaxslides,nofooterlogo]{pureminimalistic}
\usefonttheme{professionalfonts}
\usepackage[UTF8,fontset=none]{ctex}
\usepackage{graphicx}
\usepackage{listings}
\usepackage{appendixnumberbeamer}

\setsansfont{Fira Sans}
\setmonofont{Fira Mono}
\setCJKmainfont{Noto Sans CJK SC}
\setCJKsansfont{Noto Sans CJK SC}
\setCJKmonofont{Noto Sans CJK SC}

% 默认不引用上游示例 logo，有真实素材时再统一配置
\renewcommand{\logotitle}{}
\renewcommand{\logoheader}{}
\renewcommand{\logofooter}{}
\renewcommand{\pageword}{}

% 字号设置必须放在加载主题之后，避免被主题覆盖
\setbeamerfont{title}{size=\LARGE,shape=\upshape,series=\bfseries}
\setbeamerfont{presentation title}{shape=\upshape}
\setbeamerfont{frametitle}{size=\Large,shape=\upshape,series=\bfseries}

% 对浅色/深色都保持明确的代码对比度
\definecolor{TalkCodeBg}{HTML}{F3F4F6}
\definecolor{TalkCodeFg}{HTML}{17212B}
\definecolor{TalkCodeKeyword}{HTML}{174EA6}
\definecolor{TalkCodeComment}{HTML}{47634D}
\definecolor{TalkCodeString}{HTML}{8A3014}
\lstdefinestyle{talkcode}{
  basicstyle=\ttfamily\small\color{TalkCodeFg},
  backgroundcolor=\color{TalkCodeBg},
  keywordstyle=\bfseries\color{TalkCodeKeyword},
  commentstyle=\color{TalkCodeComment},
  stringstyle=\color{TalkCodeString},
  columns=fullflexible,
  keepspaces=true,
  showstringspaces=false,
  breaklines=true,
  tabsize=4,
  numbers=none,
  frame=none
}
\lstset{style=talkcode}
 引用。
% 编译时需要把仓库根目录加入 TEXINPUTS，详见根目录 AGENTS.md 第 2 节。

\usetheme[customfont,showmaxslides,nofooterlogo]{pureminimalistic}
\usefonttheme{professionalfonts}
\usepackage[UTF8,fontset=none]{ctex}
\usepackage{graphicx}
\usepackage{listings}
\usepackage{appendixnumberbeamer}

\setsansfont{Fira Sans}
\setmonofont{Fira Mono}
\setCJKmainfont{Noto Sans CJK SC}
\setCJKsansfont{Noto Sans CJK SC}
\setCJKmonofont{Noto Sans CJK SC}

% 默认不引用上游示例 logo，有真实素材时再统一配置
\renewcommand{\logotitle}{}
\renewcommand{\logoheader}{}
\renewcommand{\logofooter}{}
\renewcommand{\pageword}{}

% 字号设置必须放在加载主题之后，避免被主题覆盖
\setbeamerfont{title}{size=\LARGE,shape=\upshape,series=\bfseries}
\setbeamerfont{presentation title}{shape=\upshape}
\setbeamerfont{frametitle}{size=\Large,shape=\upshape,series=\bfseries}

% 对浅色/深色都保持明确的代码对比度
\definecolor{TalkCodeBg}{HTML}{F3F4F6}
\definecolor{TalkCodeFg}{HTML}{17212B}
\definecolor{TalkCodeKeyword}{HTML}{174EA6}
\definecolor{TalkCodeComment}{HTML}{47634D}
\definecolor{TalkCodeString}{HTML}{8A3014}
\lstdefinestyle{talkcode}{
  basicstyle=\ttfamily\small\color{TalkCodeFg},
  backgroundcolor=\color{TalkCodeBg},
  keywordstyle=\bfseries\color{TalkCodeKeyword},
  commentstyle=\color{TalkCodeComment},
  stringstyle=\color{TalkCodeString},
  columns=fullflexible,
  keepspaces=true,
  showstringspaces=false,
  breaklines=true,
  tabsize=4,
  numbers=none,
  frame=none
}
\lstset{style=talkcode}
 引用。
% 编译时需要把仓库根目录加入 TEXINPUTS，详见根目录 AGENTS.md 第 2 节。

\usetheme[customfont,showmaxslides,nofooterlogo]{pureminimalistic}
\usefonttheme{professionalfonts}
\usepackage[UTF8,fontset=none]{ctex}
\usepackage{graphicx}
\usepackage{listings}
\usepackage{appendixnumberbeamer}

\setsansfont{Fira Sans}
\setmonofont{Fira Mono}
\setCJKmainfont{Noto Sans CJK SC}
\setCJKsansfont{Noto Sans CJK SC}
\setCJKmonofont{Noto Sans CJK SC}

% 默认不引用上游示例 logo，有真实素材时再统一配置
\renewcommand{\logotitle}{}
\renewcommand{\logoheader}{}
\renewcommand{\logofooter}{}
\renewcommand{\pageword}{}

% 字号设置必须放在加载主题之后，避免被主题覆盖
\setbeamerfont{title}{size=\LARGE,shape=\upshape,series=\bfseries}
\setbeamerfont{presentation title}{shape=\upshape}
\setbeamerfont{frametitle}{size=\Large,shape=\upshape,series=\bfseries}

% 对浅色/深色都保持明确的代码对比度
\definecolor{TalkCodeBg}{HTML}{F3F4F6}
\definecolor{TalkCodeFg}{HTML}{17212B}
\definecolor{TalkCodeKeyword}{HTML}{174EA6}
\definecolor{TalkCodeComment}{HTML}{47634D}
\definecolor{TalkCodeString}{HTML}{8A3014}
\lstdefinestyle{talkcode}{
  basicstyle=\ttfamily\small\color{TalkCodeFg},
  backgroundcolor=\color{TalkCodeBg},
  keywordstyle=\bfseries\color{TalkCodeKeyword},
  commentstyle=\color{TalkCodeComment},
  stringstyle=\color{TalkCodeString},
  columns=fullflexible,
  keepspaces=true,
  showstringspaces=false,
  breaklines=true,
  tabsize=4,
  numbers=none,
  frame=none
}
\lstset{style=talkcode}
